\documentclass[11pt]{article}
\usepackage[margin=1in]{geometry}
\usepackage{fontspec}
\usepackage{unicode-math}
\usepackage{amsmath}
\usepackage{booktabs,longtable,array}
\usepackage{enumitem}
\usepackage{fancyvrb}
\usepackage[hidelinks]{hyperref}
\usepackage{microtype}
\setmainfont{Latin Modern Roman}
\setmathfont{Latin Modern Math}
\setlength{\parindent}{0pt}
\setlength{\parskip}{0.6em}
\begin{document}
\section{Continual Learning Benchmark Protocol Specification}

Version: 1.0.0

\subsection{Overview}

This document specifies a standardized protocol for evaluating continual learning algorithms. The protocol is designed to be framework-agnostic and provides reproducible, comparable evaluations across different methods.

\subsection{1. Definitions}

\subsubsection{1.1 Task Sequence}

A continual learning benchmark consists of a sequence of T tasks: ${\mathrm{Task}_0, \mathrm{Task}_1, \ldots, \mathrm{Task}_{T-1}}$.

Each task $\mathrm{Task}_t$ contains:

\begin{itemize}[leftmargin=*]
\item A training set $D_t^{\mathrm{train}} = {(x_i, y_i)}$
\item A test set $D_t^{\mathrm{test}} = {(x_i, y_i)}$
\item A set of target classes $C_t$
\end{itemize}

Tasks are presented sequentially. When training on $\mathrm{Task}_t$, only $D_t^{\mathrm{train}}$ is available (unless explicitly using replay or other memory mechanisms).

\subsubsection{1.2 Model Interface}

A model $M$ must implement two operations:

\begin{Verbatim}[fontsize=\small]
predict(x) -> y_hat
    Input: x with shape (batch_size, ...)
    Output: y_hat with shape (batch_size, num_classes)

train_on_task(task_id, train_data, epochs) -> metrics
    Input: task identifier, training data, number of epochs
    Output: dictionary of training metrics
\end{Verbatim}

\subsection{2. Evaluation Protocol}

\subsubsection{2.1 Training Phase}

For each task $t$ in sequence $0, 1, \ldots, T-1$:

\begin{enumerate}[leftmargin=*]
\item \textbf{Train}: Call $model.train_on_task(t, \mathrm{Task}_t, epochs)$
\item \textbf{Evaluate}: After training, evaluate on ALL tasks \texttt{0, 1, ..., t}
\item \textbf{Record}: Store accuracy for each evaluated task
\end{enumerate}

\subsubsection{2.2 Accuracy Matrix}

The accuracy matrix $A$ has shape \texttt{(T, T)} where:

\[
A[i, j] = \text{accuracy on } \mathrm{Task}_j \text{ after training through } \mathrm{Task}_i
\]

The matrix is lower triangular (only $A[i, j]$ where $j \le  i$ is defined).

\textbf{Evaluation procedure for $A[i, j]$:}

\begin{enumerate}[leftmargin=*]
\item For each sample \texttt{(x, y)} in $D_j^{\mathrm{test}}$:
\item Compute $y_hat = \operatorname{argmax}(model.predict(x))$
\item $A[i, j] = \operatorname{mean}(y_hat == \operatorname{argmax}(y))$
\end{enumerate}

\subsubsection{2.3 Multiple Runs}

For statistical validity, run the protocol multiple times (recommended: 5+ runs) with different random seeds. Report mean and standard deviation of all metrics.

\subsection{3. Metrics}

\subsubsection{3.1 Average Accuracy}

Final average accuracy across all tasks:

\[
Average Accuracy = (1/T) * sum_{t=0}^{T-1} A[T-1, t]
\]

\subsubsection{3.2 Forgetting}

Per-task forgetting measures the drop from peak accuracy to final accuracy:

\[
F_t = max_{i \ge  t} A[i, t] - A[T-1, t]
\]

Average forgetting (excluding the last task, which has no opportunity to be forgotten):

\[
Average Forgetting = (1/(T-1)) * sum_{t=0}^{T-2} F_t
\]

\subsubsection{3.3 Backward Transfer (BWT)}

Backward transfer measures how learning new tasks affects performance on previous tasks:

\[
BWT = (1/(T-1)) * sum_{t=0}^{T-2} (A[T-1, t] - A[t, t])
\]

\begin{itemize}[leftmargin=*]
\item \texttt{BWT < 0}: Negative transfer (forgetting)
\item \texttt{BWT > 0}: Positive transfer (learning new tasks improves old ones)
\item $BWT = 0$: No transfer effect
\end{itemize}

\subsubsection{3.4 Forward Transfer (FWT)}

Forward transfer measures how knowledge from previous tasks helps learning new tasks:

\[
FWT = (1/(T-1)) * sum_{t=1}^{T-1} (A[t-1, t] - B[t])
\]

Where $B[t]$ is the baseline accuracy on task $t$ (e.g., random initialization performance).

\begin{itemize}[leftmargin=*]
\item \texttt{FWT > 0}: Positive forward transfer
\item \texttt{FWT < 0}: Negative forward transfer (interference)
\end{itemize}

\subsubsection{3.5 Learning Curve Area}

For each task $t$, the learning curve tracks $A[i, t]$ for $i \ge  t$. The area under this curve provides a comprehensive view of both initial learning and retention.

\subsection{4. Standard Datasets}

\subsubsection{4.1 Split-MNIST}

\begin{itemize}[leftmargin=*]
\item \textbf{Tasks}: 5
\item \textbf{Classes}: 10 (digits 0-9)
\item \textbf{Task composition}: ${(0,1), (2,3), (4,5), (6,7), (8,9)}$
\item \textbf{Input shape}: (784,) or (28, 28)
\item \textbf{Labels}: One-hot encoded (10,)
\end{itemize}

\subsubsection{4.2 Permuted-MNIST}

\begin{itemize}[leftmargin=*]
\item \textbf{Tasks}: Configurable (typically 10-20)
\item \textbf{Classes}: 10 (same digits each task)
\item \textbf{Task composition}: Same classes, different pixel permutations
\item \textbf{Input shape}: (784,)
\item \textbf{Labels}: One-hot encoded (10,)
\end{itemize}

\subsubsection{4.3 Split-CIFAR10}

\begin{itemize}[leftmargin=*]
\item \textbf{Tasks}: 5
\item \textbf{Classes}: 10
\item \textbf{Task composition}: ${(airplane, automobile), (bird, cat), \ldots}$
\item \textbf{Input shape}: (3072,) or (32, 32, 3)
\item \textbf{Labels}: One-hot encoded (10,)
\end{itemize}

\subsubsection{4.4 Split-CIFAR100}

\begin{itemize}[leftmargin=*]
\item \textbf{Tasks}: 10 or 20
\item \textbf{Classes}: 100
\item \textbf{Task composition}: Groups of 10 or 5 classes
\item \textbf{Input shape}: (3072,) or (32, 32, 3)
\item \textbf{Labels}: One-hot encoded (100,)
\end{itemize}

\subsection{5. Reporting Standards}

\subsubsection{5.1 Required Information}

When reporting results, include:

\begin{enumerate}[leftmargin=*]
\item \textbf{Dataset}: Name and configuration
\item \textbf{Model architecture}: Network structure, number of parameters
\item \textbf{Training}: Epochs per task, batch size, learning rate, optimizer
\item \textbf{Evaluation}: Number of runs, random seeds used
\item \textbf{Metrics}: Average Accuracy, Forgetting, BWT (with standard deviations)
\end{enumerate}

\subsubsection{5.2 Recommended Format}

\begin{Verbatim}[fontsize=\small]
Dataset: Split-MNIST (5 tasks)
Model: MLP (784-256-10)
Training: 5 epochs/task, batch=256, lr=0.01, SGD
Runs: 5 (seeds: 0, 1, 2, 3, 4)

Results:
- Average Accuracy: 0.8534 +/- 0.0123
- Average Forgetting: 0.1245 +/- 0.0089
- Backward Transfer: -0.0987 +/- 0.0056
\end{Verbatim}

\subsubsection{5.3 Visualization}

Include at minimum:

\begin{enumerate}[leftmargin=*]
\item Accuracy matrix heatmap
\item Per-task forgetting bar chart
\item Learning curves (optional but recommended)
\end{enumerate}

\subsection{6. Statistical Significance}

\subsubsection{6.1 Comparing Methods}

When comparing two methods A and B:

\begin{enumerate}[leftmargin=*]
\item \textbf{Paired t-test}: For normally distributed differences
\item \textbf{Wilcoxon signed-rank}: For non-parametric comparison
\item \textbf{Bootstrap confidence intervals}: For robust estimation
\end{enumerate}

\subsubsection{6.2 Effect Size}

Report Cohen's d for practical significance:

\[
d = (mean_A - mean_B) / pooled_std
\]

Interpretation:

\begin{itemize}[leftmargin=*]
\item \texttt{|d| < 0.2}: Negligible
\item $0.2 \le  |d| < 0.5$: Small
\item $0.5 \le  |d| < 0.8$: Medium
\item $|d| \ge  0.8$: Large
\end{itemize}

\subsection{7. Reference Baselines}

\subsubsection{7.1 Naive Fine-tuning (Lower Bound)}

\begin{itemize}[leftmargin=*]
\item Train sequentially on each task
\item No mechanism to prevent forgetting
\item Establishes lower bound for continual learning performance
\end{itemize}

\subsubsection{7.2 Joint Training (Upper Bound)}

\begin{itemize}[leftmargin=*]
\item Train on all tasks simultaneously (all data available)
\item Violates continual learning constraints
\item Establishes upper bound for achievable performance
\end{itemize}

\subsubsection{7.3 Experience Replay}

\begin{itemize}[leftmargin=*]
\item Maintain buffer of samples from previous tasks
\item Mix replay samples with current task training
\item Common practical approach
\end{itemize}

\subsubsection{7.4 EWC (Elastic Weight Consolidation)}

\begin{itemize}[leftmargin=*]
\item Compute Fisher information after each task
\item Add quadratic penalty to protect important parameters
\item Regularization-based approach
\end{itemize}

\subsection{8. Implementation Notes}

\subsubsection{8.1 Data Format}

\begin{itemize}[leftmargin=*]
\item Images: NumPy arrays, dtype float32, range [0, 1]
\item Labels: One-hot encoded NumPy arrays, dtype float32
\item Batching: Provided through iterator interface
\end{itemize}

\subsubsection{8.2 Reproducibility}

For reproducible results:

\begin{enumerate}[leftmargin=*]
\item Set numpy random seed before each run
\item Use deterministic data shuffling
\item Document framework-specific seeds (PyTorch, JAX, TensorFlow)
\end{enumerate}

\subsubsection{8.3 Evaluation Timing}

Evaluate AFTER completing training on each task, not during training epochs. This ensures the accuracy matrix reflects the model state at task boundaries.

\subsection{Appendix A: Mathematical Notation}

\begin{longtable}{p{0.28\textwidth} p{0.62\textwidth}}
\toprule
Symbol & Description \\
\midrule
\endfirsthead
\toprule
Symbol & Description \\
\midrule
\endhead
T & Total number of tasks \\
t & Task index (0-indexed) \\
A[i,j] & Accuracy on task j after training through task i \\
$F_t$ & Forgetting for task t \\
$BWT$ & Backward Transfer \\
$FWT$ & Forward Transfer \\
$D_t^{\mathrm{train}}$ & Training data for task t \\
$D_t^{\mathrm{test}}$ & Test data for task t \\
$C_t$ & Classes in task t \\
\bottomrule
\end{longtable}

\subsection{Appendix B: Changelog}

\subsubsection{Version 1.0.0}

\begin{itemize}[leftmargin=*]
\item Initial protocol specification
\item Standard metrics definitions
\item Built-in dataset specifications
\item Reporting guidelines
\end{itemize}


\end{document}
